\chapter{Output di ogni sistema su una stessa risposta}
\label{app:examples}

% Source: Project/outputs/results/*/predictions.parquet (judge texts stored by the latency and energy runs),
% GroundOne served output via `groundone verify` (calibration + dev threshold, D27). Answer rt-7784 of T1, chosen
% because every system scored it and it contains all three classes; it was not chosen for any system's performance.

Questa appendice mostra che cosa restituisce ciascun sistema per una stessa risposta di T1, in modo da collegare i numeri
del Capitolo~\ref{sec:results} a output concreti. La risposta (identificativo \texttt{rt-7784}) è di tipo data-to-text:
un modello linguistico l'ha scritta a partire dalla descrizione strutturata di un ristorante, ossia un record JSON con
nome, indirizzo, categorie, orari di apertura (da martedì a domenica, 11:00--21:00), attributi, una valutazione media di
4.5 stelle e due recensioni di clienti. È stata scelta perché tutti i sistemi l'hanno valutata e perché contiene tutte e
tre le classi, non in base alle prestazioni di un qualche sistema. Le sei frasi e le relative etichette di riferimento
sono:

\begin{enumerate}
\item \emph{Mesa Verde is a popular business located in Santa Barbara, CA.} --- supportata.
\item \emph{It offers a variety of services including Garage Door Services, Restaurants, Live/Raw Food, \dots, Vegan, and
Salad.} --- supportata.
\item \emph{The business is open seven days a week from 11 am to 9 pm.} --- contraddetta (il lunedì il locale è chiuso).
\item \emph{Mesa Verde has a rating of 4.5 stars on Google and has received multiple positive reviews.} --- non menzionata
(il record non indica la fonte della valutazione e contiene due recensioni, una delle quali negativa).
\item \emph{One customer praised the restaurant's fresh, delicious, and inventive cuisine, while another was
disappointed by the taste and price of the food.} --- supportata.
\item \emph{Despite this, many customers have had positive experiences at Mesa Verde and recommend it for its fusion and
vegan food options.} --- supportata.
\end{enumerate}

\section{Verdetti frase per frase}

La Tabella~\ref{tab:app-examples} riporta il verdetto di ciascun sistema. Per GroundOne l'etichetta è quella restituita
dal servizio, ricavata dalle probabilità calibrate con la soglia stimata sul dev set ($\tau$ = 0.935 per il 0.5B e 0.925
per il 3B), e il numero è la $P(\text{supported})$ calibrata. Per gli encoder il numero è la probabilità di supporto, con
la decisione presa a 0.5. S, C e N stanno per supportata, contraddetta e non menzionata; un trattino indica un errore di
formato.

\begin{table}[htbp]
\centering
\caption{Verdetti sulla risposta \texttt{rt-7784}. Gli errori rispetto al riferimento sono in grassetto.}
\label{tab:app-examples}
\small
\begin{adjustbox}{max width=\textwidth}
\begin{tabular}{@{}lcccccc@{}}
\toprule
Sistema & 1 & 2 & 3 & 4 & 5 & 6 \\
\midrule
Riferimento & S & S & C & N & S & S \\
\midrule
GroundOne 0.5B (restituita, $P_\mathrm{sup}$) & S 0.96 & S 0.98 & C 0.50 & N 0.05 & S 0.98 & S 0.97 \\
GroundOne 3B (restituita, $P_\mathrm{sup}$) & S 0.96 & S 0.99 & C 0.58 & N 0.21 & S 0.99 & \textbf{N 0.85} \\
B2 0.5B (token generato) & S & S & C & N & S & \textbf{--} \\
B2 3B (token generato) & S & S & C & \textbf{--} & S & S \\
0.5B non addestrato (readout) & \textbf{C} & \textbf{C} & C & \textbf{C} & \textbf{C} & \textbf{C} \\
Giudice sintetico 0.5B & S & S & \textbf{S} & \textbf{S} & S & S \\
Giudice sintetico 3B & \textbf{C} & S & \textbf{S} & \textbf{S} & \textbf{N} & S \\
Giudice verboso 0.5B & \textbf{--} & \textbf{--} & \textbf{--} & \textbf{--} & \textbf{--} & \textbf{--} \\
Giudice verboso 3B & \textbf{--} & \textbf{--} & \textbf{--} & \textbf{--} & \textbf{--} & \textbf{--} \\
Giudice a parità di output 3B, logit & \textbf{C} & S & \textbf{S} & \textbf{S} & \textbf{N} & S \\
Giudice a parità di output 3B, probabilità espresse a parole & S & \textbf{C} & \textbf{N} & \textbf{S} & \textbf{C} & \textbf{N} \\
MiniCheck-RoBERTa ($P_\mathrm{sup}$) & 0.78 & 0.90 & 0.09 & 0.03 & \textbf{0.43} & \textbf{0.03} \\
HHEM-2.1-open ($P_\mathrm{sup}$) & 0.75 & 0.95 & 0.01 & 0.09 & 0.73 & \textbf{0.24} \\
Rizzo-Flow 4B & S & S & \textbf{S} & \textbf{S} & S & S \\
\bottomrule
\end{tabular}
\end{adjustbox}
\end{table}

GroundOne 0.5B classifica correttamente tutte e sei le frasi. La contraddizione nella frase~3 viene colta con una
$P(\text{supported})$ calibrata di 0.50: ben sotto la soglia, ma lontana dalla certezza, ed è esattamente ciò che una
probabilità calibrata dovrebbe esprimere per una frase corretta in tutto tranne che in un giorno della settimana. Anche
GroundOne 3B coglie entrambi gli errori, ma restituisce come non menzionata la frase~6, una generalizzazione tratta dalle
due recensioni: la sua $P(\text{supported})$ di 0.85 è inferiore alla soglia elevata imposta dal tasso di base di
RAGTruth, un caso in cui un utente che preferisce meno segnalazioni abbasserebbe la soglia. A parità di pesi, il token
generato (B2) dà la lettera corretta su cinque frasi su sei, ma sulla restante produce uno spazio (0.5B) o la parola
``We'' (3B) al posto di un verdetto. Gli encoder binari separano bene i due errori, ma respingono anche le frasi~5 o~6,
che riassumono le recensioni con parole diverse; Rizzo-Flow accetta tutte le frasi. I giudici generativi falliscono nei
modi quantificati nel Capitolo~\ref{sec:results}, come mostrano gli output grezzi che seguono.

\section{Output grezzi}

\paragraph{GroundOne 0.5B (JSON restituito, abbreviato).} Un unico forward pass su 1{,}193 token di input, nessun token
generato, 0.47 secondi (2.6 secondi per il 3B); nel JSON le frasi sono numerate a partire da zero:
{\footnotesize
\begin{verbatim}
{"design": "b", "threshold": 0.935, "n_decisions": 6,
 "usage": {"input_tokens": 1193, "output_tokens": 0},
 "sentences": [
  {"index": 2, "text": "The business is open seven days a week ...",
   "scored": true, "labels": {"three": "contradicted"},
   "answers": {"three": {"supported": 0.502, "contradicted": 0.462,
                         "not_mentioned": 0.036}},
   "flag_strict": true, "flag_lenient": true},
  {"index": 3, "text": "Mesa Verde has a rating of 4.5 stars ...",
   "scored": true, "labels": {"three": "not_mentioned"},
   "answers": {"three": {"supported": 0.051, "contradicted": 0.003,
                         "not_mentioned": 0.946}},
   "flag_strict": true, "flag_lenient": false}, ...]}
\end{verbatim}
}

\paragraph{Giudice sintetico 3B.} Una riga per frase, come richiesto, ma con la lettera sbagliata in quattro frasi:
{\footnotesize
\begin{verbatim}
1: B
2: A
3: A
4: A
5: C
6: A
\end{verbatim}
}

\paragraph{Giudice verboso 3B (primi tre verdetti).} Il ragionamento riformula la frase e cita il documento a
sostegno anche della frase contraddetta e di quella non supportata, e tutti i verdetti sono ``A''. I verdetti compaiono
su una riga a sé anziché sulla stessa riga richiesta dall'istruzione, per cui il protocollo li conta come errori di
formato:
{\footnotesize
\begin{verbatim}
1: Mesa Verde is located in Santa Barbara, CA, as stated in the document.
Verdict: A

2: Mesa Verde offers a variety of services including Garage Door Services,
   [...], as mentioned in the categories provided.
Verdict: A

3: The business is open seven days a week from 11 am to 9 pm, as indicated
   by the hours provided.
Verdict: A
\end{verbatim}
}

\paragraph{Giudice verboso 0.5B.} Ricopia le sei frasi, numerate, senza alcun verdetto.

\paragraph{Giudice a parità di output 3B, probabilità espresse a parole.} Le probabilità sono tutte one-hot e, per le
frasi~3 e~6, l'etichetta dichiarata (C) contraddice le probabilità dichiarate, che assegnano tutta la massa ad A:
{\footnotesize
\begin{verbatim}
{"1": {"label": "A", "p": [1.0, 0.0, 0.0]},
 "2": {"label": "B", "p": [0.0, 1.0, 0.0]},
 "3": {"label": "C", "p": [1.0, 0.0, 0.0]},
 "4": {"label": "A", "p": [1.0, 0.0, 0.0]},
 "5": {"label": "B", "p": [0.0, 1.0, 0.0]},
 "6": {"label": "C", "p": [1.0, 0.0, 0.0]}}
\end{verbatim}
}

\paragraph{Giudice a parità di output 0.5B, probabilità espresse a parole.} Lo stesso vettore per ogni frase:
{\footnotesize
\begin{verbatim}
{"1": {"label": "A", "p": [0.3, 0.3, 0.3]},
 "2": {"label": "A", "p": [0.3, 0.3, 0.3]},
 ...
 "6": {"label": "A", "p": [0.3, 0.3, 0.3]}}
\end{verbatim}
}
